\subsection{A local analytic obstruction to a cusp-only Picard argument}
\label{subsec:local-cusp-picard-obstruction}

The relative Picard space of the cusp family cannot by itself force the
normalised central restriction to vanish.  The obstruction is visible
directly in the toric model of \cref{thm:candidate-cusp-geometry}.  We give
the complete factor-of-automorphy calculation, including the descent
cocycle and the sign of the generic Chern class.

Retain the cusp notation of
\cref{prop:period-input-audit,thm:candidate-cusp-geometry}:
\[
 B_0=\begin{pmatrix}0&1\\-1&0\end{pmatrix},
 \qquad
 Z(s)=sB_0+C(t),
 \qquad
 C(t)=\begin{pmatrix}6\mu(t)&h(t)\\ b(t)-h(t)&\mu(t)\end{pmatrix},
 \qquad t=e^{2\pi i s},
\tag{5.17t}
\]
\[
 c_{\lambda}(t)=\exp(2\pi iC(t)\lambda),
 \qquad
 \exp(2\pi iZ(s)\lambda)=t^{B_0\lambda}c_\lambda(t)
 \quad(\lambda\in\bar\Lambda\simeq\mathbb Z^2),
\tag{5.17u}
\]
and
\[
 \Phi_\lambda(x,t)=(t^{B_0\lambda}x,t),
 \qquad
 \Psi_\lambda(p)=(c_\lambda(t(p)),1)\cdot\Phi_\lambda(p),
 \qquad
 \Psi_\lambda\Psi_\mu=\Psi_{\lambda+\mu}.
\tag{5.17v}
\]
Write \(E_v\), \(v\in\mathbb Z^2\), for the toric prime divisors of
\(Y_\epsilon\), so that \(E_0\simeq\mathrm{dP}_6\), and let
\[
 \pi:Y_\epsilon\longrightarrow N_0=Y_\epsilon/\bar\Lambda,
 \qquad
 \nu=\pi|_{E_0}:E_0\longrightarrow W
\]
be the quotient and the normalization of the central fibre.
Put
\[
 N'=\mathbb Z^2\oplus\mathbb Z,
 \qquad
 M'=\operatorname{Hom}(N',\mathbb Z).
\]
For \(m=(\ell,k)\in M'=\mathbb Z^2\oplus\mathbb Z\), fix the character
convention
\[
 \chi^m(x,t)=x^\ell t^k,
 \qquad
 \operatorname{div}(\chi^m)=
 \sum_{v\in\mathbb Z^2}\langle m,(v,1)\rangle E_v.
\tag{5.17w}
\]

\begin{proposition}[an explicit nonzero local cusp restriction]
\label{prop:local-cusp-picard-obstruction}
\Verified\ Let
\[
 H=\begin{pmatrix}1&0\\0&-6\end{pmatrix},
 \qquad
 P=B_0^{\mathsf t}H=\begin{pmatrix}0&6\\1&0\end{pmatrix},
 \qquad
 Q(v)=v^{\mathsf t}Hv=v_1^2-6v_2^2.
\tag{5.17x}
\]
Then
\[
 \mathcal D_Q=\sum_{v\in\mathbb Z^2}Q(v)E_v
\tag{5.17y}
\]
is a locally finite Cartier divisor on \(Y_\epsilon\).
The line bundle \(\mathcal O_{Y_\epsilon}(\mathcal D_Q)\) has an explicit
\((\bar\Lambda,\Psi)\)-equivariant structure and descends to a holomorphic
line bundle \(\mathcal L_Q\) on \(N_0\).  If
\[
 (a_1,a_2,a_3,a_4,a_5,a_6)
 =(e_1,e_2,e_2-e_1,-e_1,-e_2,e_1-e_2),
 \qquad C_i=E_0\cap E_{a_i},
\]
then
\[
 \nu^*(\mathcal L_Q|_W)
 \simeq
 \mathcal O_{E_0}\!\left(\sum_{i=1}^6Q(a_i)C_i\right),
\qquad
 \left(\deg_{C_i}\nu^*(\mathcal L_Q|_W)\right)_{i=1}^6
 =(-12,2,0,-12,2,0).
\tag{5.17z}
\]
In particular \(\nu^*(\mathcal L_Q|_W)\) is nontrivial.

On every nonzero fibre
\[
 F_t=\mathbb C^2/(Z(s)\mathbb Z^2+\mathbb Z^2),
\]
the first Chern class of \(\mathcal L_Q|_{F_t}\), in the ordered lattice
basis
\((\widehat\gamma,\widehat u,\widehat w,\widehat\delta)\) indexing the
columns of \([Z\mid I_2]\), is exactly
\[
 c_1(\mathcal L_Q|_{F_t})=-2q,
 \qquad
 q=u\wedge w+6\gamma\wedge\delta.
\tag{5.17aa}
\]
\end{proposition}

\begin{proof}
\emph{Local finiteness and Cartier status.}
The fan of \(Y\) consists of the cones over the cells of the locally finite
\(A_2\)-triangulation at height one.  If
\(U_\sigma\simeq\mathbb C^3\) is the affine chart belonging to a triangle
\(\sigma=\{v_0,v_1,v_2\}\), then \(U_\sigma\) meets precisely
\(E_{v_0},E_{v_1},E_{v_2}\).  Indeed, by the orbit--cone
correspondence, \(U_\sigma\cap E_v\ne\varnothing\) exactly when
\(\mathbb R_{\geq0}(v,1)\) is a face of the cone over \(\sigma\).
Thus \(\{E_v\}_{v\in\mathbb Z^2}\) is locally finite and
\(\mathcal D_Q\) is a well-defined locally finite Weil divisor.

Every triangle of the fan is unimodular, as proved in
\cref{thm:candidate-cusp-geometry}.  If \(z_0,z_1,z_2\) are the
characters dual to
\((v_0,1),(v_1,1),(v_2,1)\), then
\[
 \mathcal D_Q|_{U_\sigma}
 =\operatorname{div}\!
   \left(z_0^{Q(v_0)}z_1^{Q(v_1)}z_2^{Q(v_2)}\right).
\tag{5.17ab}
\]
Negative exponents merely make the displayed monomial meromorphic, as
permitted for a local equation of a Cartier divisor.  Hence
\(\mathcal D_Q\) is Cartier on every chart and therefore on
\(Y_\epsilon\).

\emph{The signed divisor pullback identity.}
For \(\lambda\in\bar\Lambda=\mathbb Z^2\), set
\[
 b_\lambda=B_0\lambda,
 \qquad
 \ell_\lambda=2Hb_\lambda,
 \qquad
 k_\lambda=Q(b_\lambda),
 \qquad
 m_\lambda=(\ell_\lambda,k_\lambda)\in M'.
\tag{5.17ac}
\]
The fan automorphism sends \(E_v\) to \(E_{v+b_\lambda}\).
Consequently \(\Phi_\lambda^*E_w=E_{w-b_\lambda}\), and
\begin{align*}
 \Phi_\lambda^*\mathcal D_Q-\mathcal D_Q
 &=\sum_{v\in\mathbb Z^2}
   \bigl(Q(v+b_\lambda)-Q(v)\bigr)E_v\\
 &=\sum_{v\in\mathbb Z^2}
   \bigl(2v^{\mathsf t}Hb_\lambda+
         b_\lambda^{\mathsf t}Hb_\lambda\bigr)E_v
 =\operatorname{div}(\chi^{m_\lambda}).
\end{align*}
The \(t\)-dependent horizontal torus translation
\(p\mapsto(c_\lambda(t(p)),1)\cdot p\) preserves every
torus-invariant prime divisor \(E_v\).  Hence
\[
 \boxed{\ 
 \Psi_\lambda^*\mathcal D_Q
 =\mathcal D_Q+\operatorname{div}(\chi^{m_\lambda})\ }.
\tag{5.17ad}
\]

\emph{The complete \(\Psi\)-cocycle.}
The matrices in (5.17t) and (5.17x) give
\[
 PC(t)=
 \begin{pmatrix}
 6(b(t)-h(t))&6\mu(t)\\
 6\mu(t)&h(t)
 \end{pmatrix},
\tag{5.17ae}
\]
which is symmetric.  Thus
\[
 A(t):=2PC(t)
\tag{5.17af}
\]
is a symmetric holomorphic matrix on the full cusp disc.  Define the
holomorphic units \(h_\lambda\) and meromorphic functions \(f_\lambda\) by
\[
 h_\lambda(t)=
 \exp\!\left(\pi i\lambda^{\mathsf t}A(t)\lambda\right),
 \qquad
 f_\lambda=h_\lambda(t)\chi^{m_\lambda}.
\tag{5.17ag}
\]
Equation (5.17ad) gives
\[
 \operatorname{div}(f_\lambda)
 =\Psi_\lambda^*\mathcal D_Q-\mathcal D_Q.
\tag{5.17ah}
\]

The cocharacter map underlying \(\Phi_\mu\) is
\[
 \varphi_\mu(y,y_3)=(y+y_3b_\mu,y_3),
\]
so its pullback on characters is
\[
 \varphi_\mu^*(\ell,k)
 =(\ell,k+\ell^{\mathsf t}b_\mu).
\tag{5.17ai}
\]
The quadratic identity for \(Q\) gives
\[
 m_{\lambda+\mu}=m_\mu+\varphi_\mu^*m_\lambda.
\tag{5.17aj}
\]
Indeed, the first coordinates add, while the last coordinate on the
right is
\[
 Q(b_\mu)+Q(b_\lambda)+2b_\lambda^{\mathsf t}Hb_\mu
 =Q(b_\lambda+b_\mu).
\]
Using (5.17u), (5.17v), and (5.17ai), one obtains
\begin{align*}
 \Psi_\mu^*\chi^{m_\lambda}
 &=\chi^{m_\lambda}(c_\mu(t),1)
   \chi^{\varphi_\mu^*m_\lambda}\\
 &=\exp\!\left(2\pi i\ell_\lambda^{\mathsf t}C(t)\mu\right)
   \chi^{\varphi_\mu^*m_\lambda}\\
 &=\exp\!\left(2\pi i\lambda^{\mathsf t}A(t)\mu\right)
   \chi^{\varphi_\mu^*m_\lambda}.
\tag{5.17ak}
\end{align*}
The last equality is exact because
\[
 \ell_\lambda^{\mathsf t}C\mu
 =(2HB_0\lambda)^{\mathsf t}C\mu
 =2\lambda^{\mathsf t}B_0^{\mathsf t}HC\mu
 =2\lambda^{\mathsf t}PC\mu
 =\lambda^{\mathsf t}A\mu.
\]
Since \(A=A^{\mathsf t}\), (5.17ag), (5.17aj), and (5.17ak) imply
\begin{align*}
 f_\mu\,\Psi_\mu^*f_\lambda
 &=h_\mu h_\lambda
   \exp(2\pi i\lambda^{\mathsf t}A\mu)
   \chi^{m_\mu+\varphi_\mu^*m_\lambda}\\
 &=\exp\!\left(
   \pi i(\lambda+\mu)^{\mathsf t}A(\lambda+\mu)\right)
   \chi^{m_{\lambda+\mu}}
 =f_{\lambda+\mu}.
\tag{5.17al}
\end{align*}

Use the divisor-sheaf convention
\[
 \mathcal O(D)(U)=
 \{g\in\mathcal M_Y(U):
   \operatorname{div}(g)+D|_U\geq0\}.
\]
If \(D'=D+\operatorname{div}(f)\), multiplication by \(f\) maps
\(\mathcal O(D')\) isomorphically to \(\mathcal O(D)\).  Thus
(5.17ad) and (5.17ag) define
\[
 \theta_\lambda:
 \Psi_\lambda^*\mathcal O(\mathcal D_Q)
 \xrightarrow{\ \times f_\lambda\ }
 \mathcal O(\mathcal D_Q).
\tag{5.17am}
\]
Because
\((\Psi_\lambda\Psi_\mu)^*=\Psi_\mu^*\Psi_\lambda^*\),
equation (5.17al) is exactly
\[
 \theta_{\lambda+\mu}
 =\theta_\mu\circ\Psi_\mu^*\theta_\lambda.
\tag{5.17an}
\]
This proves a complete equivariance datum, not merely invariance of the
divisor class.

\emph{Analytic descent.}
By \cref{thm:candidate-cusp-geometry}, the action on \(Y_\epsilon\) is
free and properly discontinuous.  Every point therefore has an open
neighbourhood \(U\) whose distinct translates are disjoint and on which
\(\pi|_U\) is a biholomorphism onto \(\pi(U)\).  Transport
\(\mathcal O(\mathcal D_Q)|_U\) to \(\pi(U)\), and on overlaps arising
from translates glue by (5.17am).  Equation (5.17an) is precisely the
triple-overlap condition.  The resulting holomorphic line bundle
\(\mathcal L_Q\) on \(N_0\) satisfies
\[
 \pi^*\mathcal L_Q
 \simeq\mathcal O_{Y_\epsilon}(\mathcal D_Q).
\tag{5.17ao}
\]

\emph{The normalized central restriction.}
Let \(j_0:E_0\hookrightarrow Y_\epsilon\).  Since
\(\nu=\pi\circ j_0\), (5.17ao) gives
\[
 \nu^*(\mathcal L_Q|_W)
 =j_0^*\pi^*\mathcal L_Q
 \simeq\mathcal O_{E_0}(\mathcal D_Q|_{E_0}).
\tag{5.17ap}
\]
The only \(E_v\), \(v\ne0\), meeting \(E_0\) are those indexed by the
six \(a_i\) in the statement, and the intersections are the six toric
\((-1)\)-curves \(C_i\).  Since \(Q(0)=0\), there is no
\(E_0|_{E_0}\) term, so
\[
 \mathcal D_Q|_{E_0}
 =\sum_{i=1}^6Q(a_i)C_i,
 \qquad
 (Q(a_1),\ldots,Q(a_6))=(1,-6,-5,1,-6,-5).
\tag{5.17aq}
\]
The hexagonal intersection table is
\[
 C_i^2=-1,
 \qquad
 C_i\cdot C_{i\pm1}=1,
 \qquad
 C_i\cdot C_j=0
 \quad(j\notin\{i,i\pm1\}),
\tag{5.17ar}
\]
with indices modulo six.  Therefore
\[
 \deg_{C_i}\mathcal O_{E_0}
 \!\left(\sum_jQ(a_j)C_j\right)
 =-Q(a_i)+Q(a_{i-1})+Q(a_{i+1}),
\]
and direct substitution gives
\[
 \begin{array}{c|rrrrrr}
 i&1&2&3&4&5&6\\ \hline
 -Q(a_i)+Q(a_{i-1})+Q(a_{i+1})
 &-12&2&0&-12&2&0.
 \end{array}
\tag{5.17as}
\]
The degree on \(C_1=C_{e_1}\) is \(-12\), proving nontriviality.

\emph{The generic-fibre factor and its exact Chern class.}
Fix \(t\ne0\), choose \(s\) with \(t=e^{2\pi i s}\), and write
\(x=e^{2\pi i\zeta}\), \(\zeta\in\mathbb C^2\).
Equation (5.17u) identifies the deck transformation indexed by
\(\lambda\in\mathbb Z^2\) with
\[
 \zeta\longmapsto\zeta+Z(s)\lambda
\]
modulo the horizontal translations
\(\zeta\mapsto\zeta+n\), \(n\in\mathbb Z^2\).  The matrix identities
\[
 HB_0=P^{\mathsf t},
 \qquad
 B_0^{\mathsf t}HB_0=PB_0
\tag{5.17at}
\]
give
\[
 \chi^{m_\lambda}(x,t)
 =x^{2P^{\mathsf t}\lambda}
  t^{\lambda^{\mathsf t}PB_0\lambda},
\qquad
 h_\lambda(t)
 =\exp(2\pi i\lambda^{\mathsf t}PC(t)\lambda).
\]
Since \(Z=sB_0+C\), it follows that
\[
 f_\lambda(\zeta)
 =\exp\!\left(
 2\pi i\lambda^{\mathsf t}PZ\lambda+
 4\pi i\lambda^{\mathsf t}P\zeta\right).
\tag{5.17au}
\]
The matrix \(PZ=sPB_0+PC\) is symmetric:
\[
 PB_0=\operatorname{diag}(-6,1),
\]
and \(PC\) is symmetric by (5.17ae).

Under the descent convention (5.17am), a pulled-back local section
written as \(\varphi(\zeta)\) in the canonical frame satisfies
\[
 \varphi(\zeta+Z\lambda)
 =f_\lambda(\zeta)^{-1}\varphi(\zeta).
\]
Thus the factor of automorphy for the lattice translation
\(\zeta\mapsto\zeta+Z\lambda+n\) is
\[
 a_{(\lambda,n)}(\zeta)
 =\exp(2\pi i g_{(\lambda,n)}(\zeta)),
\qquad
 g_{(\lambda,n)}(\zeta)
 =-\lambda^{\mathsf t}PZ\lambda
  -2\lambda^{\mathsf t}P\zeta.
\tag{5.17av}
\]
This is independent of \(n\) as a factor of automorphy: replacing
\(\zeta\) by \(\zeta+n\) changes \(g\) by the integer
\(-2\lambda^{\mathsf t}Pn\).

For \(l=(\lambda,n)\), \(l'=(\mu,m)\), the exponential connecting
cocycle is
\begin{align*}
 c(l,l')
 &=g_l(\zeta+Z\mu+m)+g_{l'}(\zeta)-g_{l+l'}(\zeta)\\
 &=-2\lambda^{\mathsf t}Pm\in\mathbb Z,
\tag{5.17aw}
\end{align*}
where the symmetry of \(PZ\) cancels the two mixed \(Z\)-terms.
The corresponding alternating integral form is
\[
 E(l,l')=c(l,l')-c(l',l)
 =-2\lambda^{\mathsf t}Pm+2\mu^{\mathsf t}Pn.
\tag{5.17ax}
\]
This is \(c_1(\mathcal L_Q|_{F_t})\).  In the ordered column basis of
\([Z\mid I_2]\), the cross-block of the class
\(q=u\wedge w+6\gamma\wedge\delta\) from
\((\widehat\gamma,\widehat u)\) to
\((\widehat w,\widehat\delta)\) is
\[
 \begin{pmatrix}
 q(\widehat\gamma,\widehat w)&
 q(\widehat\gamma,\widehat\delta)\\
 q(\widehat u,\widehat w)&
 q(\widehat u,\widehat\delta)
 \end{pmatrix}
 =\begin{pmatrix}0&6\\1&0\end{pmatrix}=P.
\tag{5.17ay}
\]
Hence (5.17ax) is exactly \(-2q\), as asserted.
\end{proof}

\begin{remark}[the exact obstruction to the purely analytic route]
\label{rem:local-cusp-picard-obstruction}
The proposition proves that the local homomorphism
\[
 \operatorname{Pic}(N_0)\longrightarrow\operatorname{Pic}(E_0),
 \qquad
 \mathcal L\longmapsto\nu^*(\mathcal L|_W),
\]
is not zero.  Thus the toric cusp geometry, conductor descent on \(W\),
or local relative-Picard representability cannot alone prove the desired
global vanishing
\[
 \nu^*i^*:\operatorname{Pic}(X)\longrightarrow
 \operatorname{Pic}(E_0).
\]
There is no contradiction with
\cref{prop:leray-edge-restriction,prop:direct-boundary-picard}: the local
bundle above has nonzero generic class \(-2q\), whereas the two global
proofs use the order-three and order-four fillings and their common
boundary relation to exclude every nonzero multiple of \(q\).

More generally, Bingener proves that if a holomorphic map is proper,
flat, and cohomologically flat in dimension zero, then its relative
Picard functor is represented by a \emph{not necessarily separated}
complex space
\cite[Introduction, p.~318; Satz~6.2, pp.~336--337]{Bingener1980}.
The same introduction explains that the exponential-sequence quotient
works directly under topological local triviality and cohomological
flatness in dimension one, but that passage to the quotient fails in
general.  The cusp map of \cref{thm:candidate-cusp-geometry} is not
topologically locally trivial across \(0\).  Bingener's theorem therefore
supplies representability but neither separatedness nor the
component-specialization assertion needed here.

Consequently a purely analytic completion must add global compatibility
with both finite fillings and must control the boundary support function;
neither relative-Picard representability nor conductor descent supplies
those statements.  The later calculation in
\cref{prop:analytic-ah-boundary,lem:analytic-support-specialization,%
thm:analytic-all-picard} now supplies them: it first forces the generic
Appell--Humbert form to vanish using the two finite fillings, and then
proves that its mixed block is the polarization of the three normalized
central degrees.  Thus this proposition remains the sharp cusp-local
obstruction, while the global purely analytic route is complete.
\end{remark}
